\documentclass[11pt]{article}
\usepackage[letterpaper,margin=0.68in,headheight=15pt]{geometry}
\usepackage[T1]{fontenc}
\usepackage{lmodern}
\usepackage{amsmath,amssymb}
\usepackage{array,booktabs,enumitem,multicol}
\usepackage{xcolor}
\usepackage{tikz}
\usepackage{fancyhdr}
\usepackage{microtype}
\setlength{\parindent}{0pt}
\setlength{\parskip}{5pt}
\definecolor{olive}{HTML}{555B35}
\definecolor{gold}{HTML}{B18A22}
\definecolor{soft}{HTML}{F4F1DE}
\definecolor{bluegray}{HTML}{E8EEF2}
\pagestyle{fancy}
\fancyhf{}
\fancyhead[L]{\textcolor{olive}{\small MATH\& 141: Precalculus I}}
\fancyhead[R]{\textcolor{gold}{\small Fall 2026}}
\fancyfoot[C]{\thepage}
\renewcommand{\headrulewidth}{0.45pt}
\renewcommand{\footrulewidth}{0pt}
\setlist[enumerate]{leftmargin=*,itemsep=0.08in,topsep=0.05in}
\newcommand{\assessmenttitle}[2]{%
  {\Large\bfseries Module #1: #2}\par\vspace{3pt}
}
\newcommand{\studentline}{%
\noindent\textbf{Name:}\hspace{2.6in}\textbf{Date:}\hspace{1.15in}\par\vspace{7pt}}
\newcommand{\directions}[1]{\textit{#1}\par\vspace{4pt}}
\newcommand{\answerblank}[1]{\vspace{#1}}
\newcommand{\blankgrid}[4]{%
\begin{center}
\begin{tikzpicture}[x=0.75cm,y=0.75cm]
  \draw[step=1,gray!28,very thin] (#1,#3) grid (#2,#4);
  \draw[->,thick] (#1,0)--(#2,0) node[right] {$x$};
  \draw[->,thick] (0,#3)--(0,#4) node[above] {$y$};
  \foreach \x in {#1,...,#2}{\ifnum\x=0\else\draw (\x,0.10)--(\x,-0.10) node[below=2pt,font=\scriptsize]{\x};\fi}
  \foreach \y in {#3,...,#4}{\ifnum\y=0\else\draw (0.10,\y)--(-0.10,\y) node[left=2pt,font=\scriptsize]{\y};\fi}
\end{tikzpicture}
\end{center}}
\newcommand{\smallgrid}[4]{%
\begin{tikzpicture}[x=0.50cm,y=0.50cm]
  \draw[step=1,gray!28,very thin] (#1,#3) grid (#2,#4);
  \draw[->,thick] (#1,0)--(#2,0) node[right] {$x$};
  \draw[->,thick] (0,#3)--(0,#4) node[above] {$y$};
\end{tikzpicture}}
\begin{document}

% MODULE 4 PREQUIZ
\assessmenttitle{4}{Rational Functions -- Prequiz}
\studentline
\directions{Four short questions. Distinguish canceled factors from denominator factors that remain.}
\begin{enumerate}
\item Divide $x^2+5x+6$ by $x+2$.
\answerblank{0.52in}
\item For
\[
f(x)=\frac{x-4}{(x-1)(x+2)},
\]
find the vertical asymptotes and the $x$-intercept.
\answerblank{0.58in}
\item For
\[
g(x)=\frac{(x+3)(x-1)}{(x+3)(x+2)},
\]
find the hole as an ordered pair.
\answerblank{0.65in}
\item A rational function has numerator degree $1$ and denominator degree $2$. What is its horizontal asymptote?
\answerblank{0.42in}
\end{enumerate}
\newpage

% MODULE 4 CHECKIN PAGE 1
\assessmenttitle{4}{Rational Functions -- Weekly Check-In}
\studentline
\directions{Three questions. Distinguish canceled factors from denominator factors that remain.}
\textbf{1. Polynomial division.}

Divide $2x^2+3x-5$ by $x-2$. Show the long division and then rewrite
\[
R(x)=\frac{2x^2+3x-5}{x-2}
\]
in quotient-remainder form.
\answerblank{2.25in}

\textbf{2. A hole and asymptotes.}

Consider
\[
f(x)=\frac{(x+2)(x-4)}{(x+2)(x-1)(x+3)}.
\]
Find the domain restrictions, the hole as an ordered pair, the vertical asymptotes, and the $x$-intercept.
\answerblank{1.35in}
\newpage

% MODULE 4 CHECKIN PAGE 2
\assessmenttitle{4}{Rational Functions -- Weekly Check-In, continued}
\studentline
\textbf{3. Intercepts and horizontal behavior.}

Consider
\[
g(x)=\frac{2x-6}{x^2+x-6}.
\]
\begin{enumerate}[label=(\alph*),leftmargin=0.35in]
\item Factor completely and state the domain restrictions.
\answerblank{0.60in}
\item Find any hole, vertical asymptote, $x$-intercept, and $y$-intercept.
\answerblank{1.20in}
\item Find the horizontal asymptote and explain how the numerator and denominator degrees determine it.
\answerblank{0.90in}
\end{enumerate}

\end{document}
